\documentclass{leveldoc}
\DocUses{foundations,deflection_plates,power_supply}
\DocFingerprint{50fc18a32d73bf35}

\begin{document}

\section{What the level asks}\label{sec:asks}

\leveldescription

Three positive ions are launched from the left edge, one straight along $+x$
and two at small angles to it. Three metal apertures stand on the launch
line. The outer apertures are grounded. Each half of the middle aperture has
its own supply. The player sets those two supplies. At the reference both are
$\val{vSupply}$. A positive middle aperture slows a positive ion as the ion
climbs the potential, and the transverse field in the openings turns the
angled rays back toward the launch line.

Each ion must enter the detector. No shot has an acceptance window, a kinetic
window or a gate. A flight ends when the ion touches a plate, leaves the
world, or enters the detector (\xsecref{foundations}{sec:board}). A plate's
obstacle is its box plus a contact shell of thickness $\val{contact}$. Each
ion's radius is $0$.

With both supplies off, the on-axis ray still follows the launch line into
the detector, and the angled rays stay apart and miss it. This document uses
the deflection plates' rectangle collocation, with each supply entered as the
potential of the electrode it sits on, and integrates the three flights by
the classical Runge--Kutta method.

\section{Setup}\label{sec:setup}

\setuptable

The speed of light is $c = \val{cLight}$. Each ion has charge $q = \val{qPart}$,
mass $m = \val{mPart}$ and radius $0$, and each is launched with kinetic
energy $T_0 = \val{T0}$ from $(\val{xLaunch}, \val{yLaunch})$. Then
$\gamma_0 = \val{gamma0}$, $\abs{\vect{p}_0} = \val{p0}$ and
$v_0 = \val{v0} = \val{vOverC}\,c$ (\xeqref{foundations}{eq:launch}). The
three directions make angles $\val{aLow}$, $\val{aAxis}$ and $\val{aHigh}$
degrees with $+x$.

The apertures stand at $x = \val{xIn}$, $x = \val{xMid}$ and $x = \val{xOut}$.
Each is a pair of plates. A plate's long direction is $y$ and its thickness
is along $x$, from $\val{xP0a}$ to $\val{xP0b}$ at the first aperture and the
same thickness at the other two. The upper plates run from $\val{yP0a}$ to
$\val{yP0b}$, and the lower plates from $\val{yP1a}$ to $\val{yP1b}$. The
opening between the inner rims is $\val{opening}$, from $\val{yRimLo}$ to
$\val{yRimHi}$. In the reference the two halves at $x = \val{xMid}$ are at
$V = \val{vSupply}$. The four outer halves are at $\val{vOuter}$. The slider
runs from $-\val{vSlider}$ to $\val{vSlider}$.

The detector, the same one for every shot, is $\val{xDetLo} \le x \le \val{xDetHi}$,
$\val{yDetLo} \le y \le \val{yDetHi}$. The world is $-1 \le x \le \val{xWorld}$,
$-1 \le y \le \val{yWorld}$. The reported time step is $h = \val{H}$. The mesh
is the deflection plates' mesh: rectangles of edge $\val{panelFine}$,
$\val{nPanels}$ of them on the upper half.

\begin{figure}[ht]
\centering
\begin{tikzpicture}
\begin{axis}[axis equal image, width=0.92\linewidth,
  xmin=-1.5, xmax=32.5, ymin=-2, ymax=22.5,
  xlabel={$x$}, ylabel={$y$}, tick label style={font=\small}]
  \draw[black!40, dashed] (-1,-1) rectangle (\rawval{xWorld},\rawval{yWorld});
  \draw[black!25] (0,0) rectangle (30,20);
  \fill[green!50!black, fill opacity=0.18]
    (\rawval{xDetLo},\rawval{yDetLo}) rectangle (\rawval{xDetHi},\rawval{yDetHi});
  \draw[green!40!black]
    (\rawval{xDetLo},\rawval{yDetLo}) rectangle (\rawval{xDetHi},\rawval{yDetHi});
  \fill[black!20] (\rawval{xP0a},\rawval{yP0a}) rectangle (\rawval{xP0b},\rawval{yP0b});
  \fill[black!20] (\rawval{xP1a},\rawval{yP1a}) rectangle (\rawval{xP1b},\rawval{yP1b});
  \fill[black!20] (\rawval{xP4a},\rawval{yP4a}) rectangle (\rawval{xP4b},\rawval{yP4b});
  \fill[black!20] (\rawval{xP5a},\rawval{yP5a}) rectangle (\rawval{xP5b},\rawval{yP5b});
  \fill[black!20] (\rawval{xP2a},\rawval{yP2a}) rectangle (\rawval{xP2b},\rawval{yP2b});
  \draw[orange!80!black]
    (\rawval{xP2a},\rawval{yP2a}) rectangle (\rawval{xP2b},\rawval{yP2b});
  \fill[black!20] (\rawval{xP3a},\rawval{yP3a}) rectangle (\rawval{xP3b},\rawval{yP3b});
  \draw[orange!80!black]
    (\rawval{xP3a},\rawval{yP3a}) rectangle (\rawval{xP3b},\rawval{yP3b});
  \addplot[orange!80!black] table {pathofflow.dat};
  \addplot[orange!80!black] table {pathoffhigh.dat};
  \addplot[orange!80!black] table {pathoffaxis.dat};
  \addplot[blue!60!black] table {pathlow.dat};
  \addplot[very thick, blue!60!black] table {pathaxis.dat};
  \addplot[blue!60!black] table {pathhigh.dat};
  \fill (\rawval{xLaunch},\rawval{yLaunch}) circle[radius=0.08];
\end{axis}
\end{tikzpicture}
\caption{The six aperture plates and the six flights. The two outlined plates
are the middle aperture, at the supply potential. The other four are
grounded. The green box is the detector. Blue paths are the reference. The
downward ray bends back up and the upward ray bends back down, and all three
meet the left side of the detector on their own side of the launch line. The
on-axis reference path is that straight line. Orange paths are the same
launches with both supplies off. The on-axis orange path lies on the same
straight line and also meets the detector. The other two stay straight and
leave through the right edge of the world, one below the detector and one
above it. Every path stays outside the plates.}
\label{fig:board}
\end{figure}

\section{Principles and model}\label{sec:model}

\begin{principles}
  \item The relativistic equation of motion \xeqref{foundations}{eq:eom}, the
    launch \xeqref{foundations}{eq:launch}, and energy in a static field
    \xeqref{foundations}{eq:energy}. There is no magnetic field, and the motion
    stays in the plane (\xsecref{foundations}{sec:plane}).
  \item A supply sets the potential of the electrode it sits on
    (\xsecref{power_supply}{sec:gap}, \physref{2.7}). A grounded electrode is
    potential $0$. The surface charge is the deflection plates' collocation
    \xeqref{deflection_plates}{eq:collocation}, with the rectangle integral
    \xeqref{deflection_plates}{eq:rectangle}.
  \item The two halves of an aperture are at the same potential, so reflection
    through the launch line is a symmetry of the plates. The focusing force is
    the panel field in the openings. The flight uses that field.
  \item The board's rules (\xsecref{foundations}{sec:board}) and the classical
    Runge--Kutta method with cubic Hermite event location
    (\xsecref{foundations}{sec:numerics}).
  \item The radiated power \xeqref{foundations}{eq:larmor}, summed by Simpson's
    rule \xeqref{foundations}{eq:simpson}. It is not put back into the force.
\end{principles}

The reference potential on each middle half is the supply value, and the
outer halves are grounded. Off the metal the field is the field of the
surface charge and of its mirror. There is no placed charge. The state
$(\vect{x}, \vect{p})$ obeys \xeqref{foundations}{eq:eom} with
$\vect{B} = \vect{0}$, from the launch until the first ending event.

\begin{assumption}{The supplies are static}\label{as:static}
Each potential equals its slider value at once. There is no resistor and no
ramp. The launch speed is $\val{vOverC}\,c$. The game uses the same static
supply (\physref{2.7}).
\end{assumption}

\begin{assumption}{No radiation reaction}\label{as:radiation}
Each ion radiates, and that energy is not taken from the flight. The largest
of the three reference integrals is a fraction $\val{radFracMax}$ of $T_0$.
\end{assumption}

\begin{assumption}{No image of the ion}\label{as:image}
The ion's own surface charge is left out, as the deflection plates leave it
out and as the game leaves it out (\physref{2.7}). The neglected force is
bounded by $q^2/d^2$, with the distance $d$ floored at $10^{-3}$, relative to
$\max(\abs{q\vect{E}}, T_0)$. The reference samples reach $\val{refImage}$,
under the engine's limit $10^{-10}$. With the supplies off, the angled rays
pass closer to the metal and those samples reach $\val{offImage}$. The
samples are the step ends and the ending point.
\end{assumption}

\begin{assumption}{Equal rectangles}\label{as:panels}
The mesh is the one stated for the deflection plates: equal rectangles of edge
$\val{panelFine}$, not the engine's graded triangles. Replacing the edge by
$\val{panelCoarse}$ moves an end state by $\val{errMesh}$. With every
electrode grounded the density is zero on every mesh, so that change does not
move the bare rays.
\end{assumption}

\begin{assumption}{The fixed step}\label{as:step}
The reported step is $h = \val{H}$, coarser than the step used for the
deflection plates. Doubling it estimates an end state by $\val{errStep}$.
Ten times the panel-size estimate $\val{errMesh}$ lies above $10^{-9}$, and
it is larger than the step estimate $\val{errStep}$, so the tolerance on a
state is that product, $\val{tolFlight}$. Between samples the ion moves at
most $\val{stepReach}$. The smallest sampled clearance is $\val{refDmin}$ on
the reference flights and $\val{offDmin}$ with the supplies off.
\end{assumption}

\section{Derivation}\label{sec:derivation}

\subsection{The supplies and the axis}\label{sec:axis}

The two supply elements sit on the two halves of the aperture at
$x = \val{xMid}$. In the reference each value is $\val{vSupply}$, so those
two electrodes are at that potential and the other four are at $\val{vOuter}$.
With the supplies omitted, all six potentials are $0$.

The two halves of every aperture are at the same potential, and the plates
are placed symmetrically about $y = \val{yLaunch}$. Reflection through that
line leaves the metal and the potentials unchanged, so the transverse field
on the line vanishes,
\begin{equation}\label{eq:axis}
  E_y(x, \val{yLaunch}) = 0.
\end{equation}
The panel solution leaves a leftover $\val{eAxis}$ at the middle of the
opening. An ion launched on the line stays on it. Only its speed changes.

The conserved diagnostic is the deflection plates' flight energy
\xeqref{deflection_plates}{eq:flight-energy}. On this line the kinetic energy
is therefore
\begin{equation}\label{eq:slow}
  T = T_0 - q\bigl(\phi - \phi_\mathrm{launch}\bigr).
\end{equation}
At the launch, $\phi = \val{phiLaunch}$. In the middle opening,
$\phi = \val{phiMid}$, so $T = \val{Tmid}$. The ion gives up a fraction
$\val{slowFrac}$ of $T_0$ and is still moving forward. On the launch line the
potential is $\val{phiIn}$ at the entrance aperture and $\val{phiOut}$ at the
exit aperture, both below the middle value. The march uses the panels.

\subsection{The focusing field and the panels}\label{sec:focus}

Above the launch line, at $(\val{xMid}, \val{yAbove})$, the panel field is
$\val{eAbove}$, toward the line. Below it, at $(\val{xMid}, \val{yBelow})$,
the field is $\val{eBelow}$, also toward the line. A positive ion is pushed
toward the launch line from either side. That is the focus. The rectangles,
the mirror and the collocation are
\xeqref{deflection_plates}{eq:rectangle} and
\xeqref{deflection_plates}{eq:collocation}. The right-hand side is the
electrode potential, because there is no placed charge. The reported mesh has
$\val{nPanels}$ upper rectangles. The residual at the centroids is
$\val{solveRes}$. On the large faces, off the centroids, the potential
differs from the electrode potential by at most $\val{surfRes}$.

The middle aperture carries charge $\val{qMid}$. The four grounded halves
together carry $\val{qOuter}$. The state is advanced by
\xeqref{foundations}{eq:rk4} with step $h = \val{H}$. Events are located as
on the deflection plates: thirty-two Hermite samples on the step, then sixty
bisections, obstacles before the world, and the world before the detector.
A face coordinate is assigned when the residual is below $10^{-6}$. The
step-doubling estimate divides the change by $15$; the panel-size estimate
is the change itself when the edge doubles.

\section{The reference solution by hand}\label{sec:hand}

\begin{procedure}
  \item\label{st:launch} \emph{The launches} \xeqref{foundations}{eq:launch}.
    $\gamma_0 = \val{gamma0}$ and $\abs{\vect{p}_0} = \val{p0}$, so
    $v_0 = \val{v0}$. The three angles are $\val{aLow}$, $\val{aAxis}$ and
    $\val{aHigh}$ degrees.
  \item\label{st:supply} \emph{The supplies.} Set both halves at
    $x = \val{xMid}$ to $\val{vSupply}$ and leave the other four electrodes
    at $\val{vOuter}$. The opening is $\val{opening}$.
  \item\label{st:panels} \emph{The panels.} Solve
    \xeqref{deflection_plates}{eq:collocation} on the mesh of
    \S\ref{sec:focus}. The centroid residual is $\val{solveRes}$. The
    off-centroid discrepancy on the large faces is $\val{surfRes}$. The
    middle aperture carries $\val{qMid}$ and the grounded halves carry
    $\val{qOuter}$. On the launch line the potential is $\val{phiLaunch}$ at
    the launch, $\val{phiIn}$ in the entrance opening, $\val{phiMid}$ in the
    middle opening and $\val{phiOut}$ in the exit opening.
  \item\label{st:slow} \emph{The slowing and the focus.}
    \eqref{eq:slow} gives $T = \val{Tmid}$ in the middle opening. The
    transverse field there is $\val{eAbove}$ above the line and
    $\val{eBelow}$ below it, each toward the line. The leftover on the line
    is $\val{eAxis}$.
  \item\label{st:march} \emph{The on-axis march.} Advance
    \xeqref{deflection_plates}{eq:flight-energy} and
    \xeqref{foundations}{eq:eom} by \xeqref{foundations}{eq:rk4} with step
    $\val{H}$. The same march, with the launch direction changed, is used for
    the angled rays. The on-axis state every five time units:
    \begin{center}\small\val{states}\end{center}
  \item\label{st:end} \emph{The arrivals.} The downward ray meets the
    \val{LowFace} face of the detector at $(\val{LowX}, \val{LowY})$, at
    $t = \val{LowT}$, with momentum $(\val{LowPx}, \val{LowPy})$, at an angle
    $\val{LowAng}$ degrees from $+x$. Its lowest point is $\val{LowYmin}$,
    above the rim $\val{yRimLo}$, and its highest is the launch line. The
    on-axis ray meets the \val{AxisFace} face at $(\val{AxisX}, \val{AxisY})$,
    at $t = \val{AxisT}$, with momentum $(\val{AxisPx}, \val{AxisPy})$, at an
    angle $\val{AxisAng}$ degrees from $+x$. The upward ray meets the
    \val{HighFace} face at $(\val{HighX}, \val{HighY})$, at
    $t = \val{HighT}$, with momentum $(\val{HighPx}, \val{HighPy})$, at an
    angle $\val{HighAng}$ degrees from $+x$. Its highest point is
    $\val{HighYmax}$, below the rim $\val{yRimHi}$. All three abscissae are
    the detector's left edge, and all three ordinates lie between
    $\val{yDetLo}$ and $\val{yDetHi}$. The runs take $\val{LowSteps}$,
    $\val{AxisSteps}$ and $\val{HighSteps}$ steps.
  \item\label{st:check} \emph{The checks.} Along the three flights, $W$ drifts
    by at most $\val{LowWdev}$, $\val{AxisWdev}$ and $\val{HighWdev}$. The
    step-doubling estimates are $\val{LowErr}$, $\val{AxisErr}$ and
    $\val{HighErr}$, and the largest of the flights is $\val{errStep}$. The
    panel-size estimate is $\val{errMesh}$. The sampled clearance is
    $\val{refDmin}$, outside the shell $\val{contact}$. The image-force ratio
    reaches $\val{refImage}$.
  \item\label{st:rad} \emph{The radiation.} Simpson's rule on steps of $h/2$
    gives $\val{LowRad}$, $\val{AxisRad}$ and $\val{HighRad}$. The largest
    fraction of $T_0$ is $\val{radFracMax}$. Halving the quadrature step again
    changes the angled integrals by the relative amount $\val{LowRadErr}$.
    These integrals are not rows of the comparison.
\end{procedure}

\begin{result}[The reference flights]
All three ions enter the detector through its left side, at
$(\val{LowX}, \val{LowY})$, $(\val{AxisX}, \val{AxisY})$ and
$(\val{HighX}, \val{HighY})$. The level is solved.
\end{result}

\section{Comparison with the game}\label{sec:comparison}

\comparisontable

Twenty-four comparisons are in the table: outcome, time and end position for
each of the three reference rays, and the same four numbers with both
supplies off. The tolerance on a state is $\val{tolFlight}$. Ten times the
panel-size estimate $\val{errMesh}$ lies above $10^{-9}$, and it is larger
than the step-doubling estimate $\val{errStep}$, so the tolerance is that
product. The tolerance on a state was fixed from these estimates before the
comparison.

The largest absolute differences are the lower reference ordinate,
$\val{cmpLowYDiff}$, and the upper reference ordinate, $\val{cmpHighYDiff}$.
Each lies above the panel-size estimate $\val{errMesh}$ and under the
tolerance. The lower and upper reference times differ by $\val{cmpLowTDiff}$
and $\val{cmpHighTDiff}$. The step estimate $\val{errStep}$ is far below
those ordinates: the panels, not the time step, set that difference. The
outcomes agree.

The three reference abscissae agree exactly, and each is the detector's left
edge. The three bare abscissae agree exactly: the angled rays at the world's
right edge, and the on-axis ray at the detector's left edge. The bare on-axis
ordinate agrees exactly, and it is the launch line. The reference on-axis
ordinate differs by $\val{cmpAxisYDiff}$, under the panel-size estimate. The
bare times and the bare angled ordinates lie under the step estimate.

\section{Other solutions}\label{sec:others}

Both supplies off, which is the empty placement in \code{placements.json},
does not solve the level. Every electrode is then grounded, the panel density
is zero, and each ray flies in a straight line. The on-axis ray still meets
the \val{OffAxisFace} face of the detector at
$(\val{OffAxisX}, \val{OffAxisY})$, at $t = \val{OffAxisT}$, with momentum
$(\val{OffAxisPx}, \val{OffAxisPy})$. That time is shorter than the reference
on-axis time $\val{AxisT}$, because the reference ion climbs the potential
in \eqref{eq:slow}. The downward ray leaves through the \val{OffLowFace} edge
of the world at $(\val{OffLowX}, \val{OffLowY})$, at $t = \val{OffLowT}$, at
an angle $\val{OffLowAng}$ degrees from $+x$. The ordinate lies below
$\val{yDetLo}$. The upward ray leaves through the \val{OffHighFace} edge at
$(\val{OffHighX}, \val{OffHighY})$, at $t = \val{OffHighT}$, at an angle
$\val{OffHighAng}$ degrees from $+x$, above $\val{yDetHi}$.

The sampled clearance with the supplies off is $\val{offDmin}$, and the
image-force ratio reaches $\val{offImage}$. The energy drift on these three
flights is $\val{OffLowWdev}$, $\val{OffAxisWdev}$ and $\val{OffHighWdev}$.
The slider's half-range is $\val{vSlider}$. This document integrates the
reference setting and the off setting.

\section{What the game computes}\label{sec:engine}

A power supply sets the potential of the tunable electrode it sits on, and
the factorization of the panel system is reused when only that potential
changes (\physref{2.7}). The outer electrodes are grounded. The surface charge
is the engine's triangle collocation, the same boundary condition as
\xeqref{deflection_plates}{eq:collocation}. This document uses the equal
rectangles of the deflection plates. The end states move by $\val{errMesh}$
when that edge doubles, and the angled reference ordinates in
\S\ref{sec:comparison} lie above that estimate and under ten times it.

The ion's image force is omitted on both sides. The bound sampled on the
reference flights reaches $\val{refImage}$, under the engine's limit
$10^{-10}$. The flight in the engine is the adaptive method of \physref{5}
and \physref{6}, verified as in \physref{7}. Our step estimate is
$\val{errStep}$. Radiation is integrated and not applied to the momentum
(\xsecref{foundations}{sec:radiation}, \physref{8}).

\section{Further reading}\label{sec:reading}

The rectangle integral, the collocation and the flight energy are derived for
the deflection plates (\xeqref{deflection_plates}{eq:rectangle},
\xeqref{deflection_plates}{eq:collocation},
\xeqref{deflection_plates}{eq:flight-energy}). A supply sets the electrode it
sits on in the power-supply level (\xsecref{power_supply}{sec:gap}). The
motion and the Runge--Kutta step are in the Foundations
(\xsecref{foundations}{sec:numerics}). The engine's supplies and electrodes
are \physref{2.7}.

\levelfileappendix

\end{document}
